\chapter{Frontend Error Boundaries}

Not every frontend error is an API error. A malformed prop or an
\texttt{undefined.map()} crashes the entire React tree to a blank screen
unless something catches it --- that's an Error Boundary.

\section{Render Errors vs. API Errors}

\begin{warnbox}[WARNING]
Error Boundaries only catch errors during rendering, lifecycle methods, and
constructors below them --- not event handlers, async code, or
\texttt{fetch}/\texttt{axios} calls. API failures still use the
try/catch + \texttt{error} state pattern from Chapter~33.
\end{warnbox}

\begin{diagrambox}[What Each Mechanism Catches]
\begin{verbatim}
 Error type                        | Caught by
 -----------------------------------|---------------------------
 Render crash (bad prop, null)      | Error Boundary
 API call fails (network, 4xx/5xx)  | try/catch + error state (Ch.33)
 Error inside onClick/onSubmit      | Local try/catch in the handler
\end{verbatim}
\end{diagrambox}

\section{Building an ErrorBoundary}

Error Boundaries must be class components --- no hook equivalent exists for
\texttt{getDerivedStateFromError}/\texttt{componentDidCatch}.

\begin{lstlisting}[language=JavaScript]
// components/ErrorBoundary.jsx
import { Component } from "react";

class ErrorBoundary extends Component {
  state = { hasError: false, error: null };

  static getDerivedStateFromError(error) {
    return { hasError: true, error }; // render phase, no side effects here
  }

  componentDidCatch(error, errorInfo) {
    console.error("ErrorBoundary caught:", error, errorInfo); // log here instead
  }

  handleReset = () => this.setState({ hasError: false, error: null });

  render() {
    if (this.state.hasError) {
      return (
        <div className="p-6 text-center">
          <h2 className="text-lg font-semibold text-red-600 mb-2">Something went wrong.</h2>
          <p className="text-sm text-gray-600 mb-4">{this.state.error?.message}</p>
          <button onClick={this.handleReset} className="px-4 py-2 bg-blue-600 text-white rounded">
            Try Again
          </button>
        </div>
      );
    }
    return this.props.children;
  }
}

export default ErrorBoundary;
\end{lstlisting}

\begin{notebox}[NOTE]
\texttt{getDerivedStateFromError} is static and side-effect free;
\texttt{componentDidCatch} runs after render and is where you log to
Sentry/LogRocket/your own endpoint.
\end{notebox}

\section{Wrapping the App}

Wrap independent sections rather than the whole app, so a crash in one
widget doesn't take down the navbar too.

\begin{lstlisting}[language=JavaScript]
// App.jsx
function App() {
  return (
    <BrowserRouter>
      <ErrorBoundary><Navbar /></ErrorBoundary>
      <ErrorBoundary>
        <Routes>
          <Route path="/tasks" element={<TaskListPage />} />
          <Route path="/profile" element={<ProfilePage />} />
        </Routes>
      </ErrorBoundary>
    </BrowserRouter>
  );
}
export default App;
\end{lstlisting}

\begin{tipbox}[TIP]
Granular boundaries (per route/widget) contain crash damage; one app-wide
boundary turns any bug into a full-page failure with no escape route.
\end{tipbox}

\section{Error Boundaries Do Not Replace API Error Handling}

Wrapping a component in an \texttt{ErrorBoundary} does nothing for a failed
\texttt{axios} call in its \texttt{useEffect} --- that rejection happens
outside React's render cycle. Every data-fetching component still needs
its own loading/error/retry handling (Chapter~33); boundaries are a last
line of defense against render crashes, not a substitute for expected
failure modes.

\whatwhyhow
{A class-based ErrorBoundary uses \texttt{getDerivedStateFromError} and \texttt{componentDidCatch} to catch render/lifecycle errors and show a fallback UI instead of a blank screen.}
{Uncaught render errors crash the whole tree with no recovery; async/API errors are invisible to boundaries and need separate handling.}
{Wrap independent sections (navbar, routes, risky widgets) in their own \texttt{<ErrorBoundary>}, and keep try/catch plus loading/error state (Chapter~33) for API calls.}
